\chapter{Dello Scorrere}

In ogni individuo, così come in ogni forma di vita, dal minuscolo insetto al leviatano, è in continuo manifestarsi una
forza inarrestabile, una volontà inconsapevole che si genera all'interno di noi senza il nostro consenso. Essa esiste
in noi prima di ogni pensiero, suscita le nostre emozioni, mette in moto le nostre reazioni;
lo fa senza riguardo per i nostri progetti o desideri e il suo fine immediato è la vita nella sua forma essenziale.

Questa forza vitale che precede ogni sé, noi la chiamiamo Fiume\index{fiume} o Scorrere ed esso è libero
nel senso che è libero di esistere secondo la propria cieca autorità. Essa è il risultato
della nostra storia naturale, del modo in cui siamo cresciuti, delle esperienze che abbiamo vissuto; è la sintesi
della civiltà che chi ha cullati, delle persone che abbiamo conosciuto e amato, delle autorità che ci hanno influenzato;
prima della nostra nascita è il risultato di tutti i poteri che la natura esercita sul mondo.
Impossibile sottrarsene, il suo solo scorrere è significato di vivere, e interromperne il flusso vuol dire l'oblio.\\

Noi percepiamo la nostra condotta come giusta, per la maggior parte, è crediamo di essere gli autori di questa giustizia;
il fine di ogni individuo è perciò, nella sua essenza più pura, agire in maniera morale. La nostra morale attuale è
l'insieme di valori\index{valore} intesi come condizioni o attributi che percepiamo come desiderabili e tali valori
noi li produciamo in tre modi per cui ne esistono tre tipi diversi.

I più comuni sono quelli adottati inconsapevolmente dal mondo in cui siamo cresciuti; questi sono generalmente accettati come veri dalla società.

Poi vengono i valori che produciamo in seguito alle nostre esperienze; questi sono in noi percepiti come più intimi, li attribuiamo
alla nostra identità. Sentiamo che esisterebbero in noi ovunque noi siamo, a prescindere dalla nostra posizione nel mondo.

Gli ultimi valori possiedono la stessa intimità dei secondi, ma essi sono il risultato del nostro ponderare, e non della semplice 
effimera esperienza. Produrli ci richiede fatica, perché ci impone di usare il pensiero e di sottoporre le nostre conclusioni al nostro tribunale
interno. Non di rado mettono in dubbio i precedenti, e attraverso di essi l'individuo sedimenta le trasformazioni più solide.\\

Tutti questi valori formano il corso del fiume, che si rende evidente ogni volta che sentiamo, in maniera intensa, le emozioni che sono suscitate tramite i nostri sensi e il nostro inconscio pensare.
Le emozioni sono, dunque, il raccordo fra il fiume e la realtà. Esse vanno considerate attentamente perché non è raro 
che il fiume si manifesti in modi che ci sorprendono; questi modi sorprendenti lo sono in quanto essi ci appaiono lontani dalla nostra morale, ma questo accade perché ci stanno fornendo nuove informazioni dal mondo e non perché la nostra identità ci sta ingannando. Non dobbiamo fare questo errore, perché il fiume non è che l'espressione di quella forza vitale che tutto abbraccia e nella quale siamo semplici navigatori. Se confondiamo quest'ultimo con noi stessi
cercheremo di forzare il flusso; proveremo a giustificare le emozioni provate come se fossero la nostra morale e per farlo useremo diversi veleni della mente, di cui parleremo, per mezzo dei quali non troveremo mai una conclusione aderente alla realtà.

